\documentclass[10pt,a4paper]{amsart}
\usepackage[margin=19mm]{geometry}
\usepackage{amsmath,amssymb,mathtools}
\usepackage[scr=rsfs,cal=euler]{mathalfa}
\usepackage{xfrac}
\usepackage[unicode,hidelinks,bookmarks=false]{hyperref}
\newcommand{\ie}{i.e.\ }
\newcommand{\cf}{cf.\ }
\newcommand{\resp}{respectively\ }
\newcommand{\Subsection}{\subsection}
\providecommand{\nobreakdash}{\nobreak\hbox{-}}
\setlength{\emergencystretch}{2em}
\multlinegap0pt
\usepackage[shortlabels]{enumitem}
\usepackage{amssymb}
\usepackage{mathtools}
\usepackage{float}
\usepackage{xcolor}
\newtheorem{theorem}{Theorem}
\title{\textbf{Stability and Bifurcation Analysis of Fractional Delay Differential Equation with a Delay-dependent Coefficient}}
\author{Pragati Dutta, Sachin Bhalekar}
\date{Source: arXiv:2605.04822v1 (CC BY 4.0)}
\begin{document}
\maketitle

\noindent\emph{Source and licence.} \textbf{Stability and Bifurcation Analysis of Fractional Delay Differential Equation with a Delay-dependent Coefficient}. Version arXiv:2605.04822v1 (CC BY 4.0), distributed by arXiv under the Creative Commons Attribution 4.0 International licence, http://creativecommons.org/licenses/by/4.0/, as recorded in the arXiv licence metadata for this exact version and shown on the abstract page https://arxiv.org/abs/2605.04822v1. Full licence notice: \texttt{licenses/arxiv-2605.04822-CC-BY-4.0.txt}.\par\smallskip

\noindent\textit{Editorial scope.} Counted unit: the article's theorem thmcase2\_unified (source0.tex lines 730--757). The article's complete original proof of it is delivered with it (source0.tex lines 758--866, one \texttt{proof} environment of 109 lines). No mathematics is written, repaired or re-derived: the statement and the proof are the article's own lines, in the article's own notation, and no retained line is substituted or re-flowed.  Retained uncounted as the source-local context the proof needs: lines 186--193.  Nothing was omitted from any retained span, so the package carries no omission record.  No retained line contains a citation, so the wrapper carries no bibliography and no bibliography entry.  No retained line contains a figure environment, an image-inclusion command or drawing code, so no image is delivered and the package declares no asset dependency.  Editorial formatting only, in addition: an AMS \texttt{amsart} preamble that restates the article's own declarations and macros used by the retained lines, namely the environments \texttt{theorem} on separate counters, together with the standard packages those lines need.  The retained environments print in this document as Theorem 1 in order of appearance, whereas the article prints the same items with the numbering of its own full document.  The complete native source of the article as distributed in the arXiv e-print for this exact version is bundled unchanged as \texttt{sources/199899/source0.tex}.
	\begin{equation}
		D^{\alpha} x(t)
		= -\gamma x(t)
		+ g\big(x(t - \tau_1)\big)
		- e^{-\gamma \tau_2}\, g\big(x(t - \tau_1 - \tau_2)\big),
		\qquad 0 < \alpha \le 1,
		\label{eq:main}
	\end{equation}

\begin{theorem}
	\label{thmcase2_unified}
	The equilibrium point $x_* = 0$ of Eq.~(\ref{eq:main}) i.e 
		\begin{equation}
		D^{\alpha} x(t)
		= -\gamma x(t)
		+ g\big(x(t - \tau_1)\big)
		- e^{-\gamma \tau_2}\, g\big(x(t - \tau_1 - \tau_2)\big),
		\qquad 0 < \alpha \le 1
	\end{equation}
	 exhibits the following stability behavior:
	
	\textbf{Delay-Independent Instability:}
	\begin{enumerate}[label=(\roman*)]
		\item If $0 < \gamma < k$, then for 
		\[
		\tau_2 > -\frac{1}{\gamma} \log\left(\frac{k - \gamma}{k}\right),
		\]
		the equilibrium is unstable for all $\tau_1 \ge 0$ and $0 < \alpha \le 1$.
		
		\item If $\gamma < 0 < k$, then for 
		\[
		\tau_2 < -\frac{1}{\gamma} \log\left(\frac{k - \gamma}{k}\right),
		\]
		the equilibrium is unstable for all $\tau_1 \ge 0$ and $0 < \alpha \le 1$.
	\end{enumerate}
	
\end{theorem}

\begin{proof}
	Consider the linearized version of  fractional delay differential equation (\ref{eq:main}) i.e 
	\begin{equation}
		D^{\alpha} x(t)
		= -\gamma x(t)
		+ k\,x(t - \tau_1)
		- k e^{-\gamma \tau_2}\, x(t - \tau_1 - \tau_2),
	\end{equation}
	where $0<\alpha\le 1$. The stability of the equilibrium $x_*=0$ is determined by the roots of the corresponding characteristic equation i.e
	\[
	\lambda^{\alpha}
	= -\gamma + k e^{-\lambda \tau_1}
	- k e^{-\gamma \tau_2} e^{-\lambda(\tau_1+\tau_2)}.
	\]
	Equivalently,
	\begin{equation}
		\lambda^{\alpha} + \gamma
		- k e^{-\lambda \tau_1}
		+ k e^{-\gamma \tau_2} e^{-\lambda(\tau_1+\tau_2)}=0.
		\label{char}
	\end{equation}
	
	\medskip

	Define
	\[
	\Delta(\lambda)
	:= \lambda^{\alpha} + \gamma
	- k e^{-\lambda \tau_1}
	+ k e^{-\gamma \tau_2} e^{-\lambda(\tau_1+\tau_2)}.
	\]
	Evaluating at $\lambda=0$, we get
	\[
	\Delta(0)
	= \gamma - k + k e^{-\gamma \tau_2}.
	\]
	
	The sign of $\Delta(0)$ plays a key role in determining the existence of positive real roots.
	
	\medskip

	
	For real $\lambda>0$, note that
	\[
	\lambda^{\alpha} \to +\infty \quad \text{as } \lambda \to +\infty,
	\]
	while the exponential terms remain bounded. Hence,
	\[
	\Delta(\lambda)\to +\infty \quad \text{as } \lambda\to+\infty.
	\]
	
	Thus, if $\Delta(0)<0$, then by continuity there exists $\lambda>0$ such that $\Delta(\lambda)=0$. This yields a real positive root, implying instability.
	
	\medskip
	\noindent

	
	We have
	\[
	\Delta(0)<0
	\quad \Longleftrightarrow \quad
	\gamma - k + k e^{-\gamma \tau_2} < 0,
	\]
	which simplifies to
	\[
	e^{-\gamma \tau_2} < \frac{k-\gamma}{k}.
	\]
	
	\medskip
	\noindent
	\textbf{Case (i): $0<\gamma<k$.}\\

	Here $\frac{k-\gamma}{k}\in(0,1)$. Taking logarithm, we get,
	\[
	-\gamma \tau_2 < \log\left(\frac{k-\gamma}{k}\right).
	\]
	
	\[
\Leftrightarrow	\tau_2 > -\frac{1}{\gamma}\log\left(\frac{k-\gamma}{k}\right).
	\]
	Thus, $\Delta(0)<0$, and hence a positive real root exists. Therefore, the equilibrium is unstable. The condition is independent of $\tau_1$, so instability holds for all $\tau_1\ge 0$.
	
	\medskip
	\noindent
	\textbf{Case (ii): $\gamma<0<k$.}
\hspace{0.2cm}	
	Now $\frac{k-\gamma}{k}>1$. Again,
	\[
	e^{-\gamma \tau_2} < \frac{k-\gamma}{k}.
	\]
	Taking logarithm, we get,
	\[
	-\gamma \tau_2 < \log\left(\frac{k-\gamma}{k}\right).
	\]

	\[
\Leftrightarrow	\tau_2 < -\frac{1}{\gamma}\log\left(\frac{k-\gamma}{k}\right).
	\]
	Thus, $\Delta(0)<0$, implying existence of a positive real root and hence instability for all $\tau_1\ge 0$.
	
	\medskip
	\noindent

	The condition for instability depends only on $\Delta(0)$, which does not involve $\tau_1$. Therefore, the instability is delay-independent with respect to $\tau_1$.
	
	\medskip
	\noindent
	Hence, in both cases, the equilibrium $x_*=0$ is unstable for all $\tau_1\ge 0$ whenever the stated condition on $\tau_2$ holds.
\end{proof}

\end{document}
